\lecture_date{15/11/2026}
\chapter{Lecture 10}

\section{Quadratic functions and quadratic forms}\label{sec:lec10-quadratics}
This lecture introduces quadratic functions in the matrix form $x^TAx+b^Tx+c$. We learn to construct this representation, differentiate it, and classify its stationary points. All variables and coefficients are real, and optimisation is \emph{unconstrained}: $x\in\mathbb R^n$.

\subsection{Matrix representation}
The single-variable quadratic $ax^2+bx+c$ has a squared term, a linear term, and a constant. In several variables, degree two also includes cross terms such as $x_1x_2$. The degree of a monomial is the sum of its exponents: $x_1x_2$ has degree two, whereas $x_1^2x_2x_3$ has degree four.

A general quadratic function of $n$ variables has the form
\[
 f(x)=\sum_{i=1}^n\alpha_i x_i^2
      +\sum_{i<j}\gamma_{ij}x_ix_j
      +\sum_{i=1}^n b_ix_i+c.
\]
There are at most $n+\binom n2+n+1$ distinct terms. For four variables this is $4+6+4+1=15$, so a matrix representation is useful.

\begin{definition}[Matrix representation]
For $x,b\in\mathbb R^n$, $A\in\mathbb R^{n\times n}$ and $c\in\mathbb R$, a quadratic function can be written as
\begin{equation}\label{eq:lec10-standard}
 \boxed{f(x)=x^TAx+b^Tx+c.}
\end{equation}
The homogeneous degree-two part $q(x)=x^TAx$ is called a \vocab{quadratic form}.
\end{definition}
The lecturer also uses ``quadratic form'' for the full expression. We distinguish the quadratic part from the full function when differentiating.

\begin{example}[Two-variable construction]\label{ex:lec10-two}
Write the function
\[
 f(x_1,x_2)=5x_1^2+7x_2^2-3x_1x_2+4x_1-2x_2+13
\]
in the form \eqref{eq:lec10-standard}.
\end{example}
Take
\[
 x=\begin{bmatrix}x_1\\x_2\end{bmatrix},\qquad
 A=\begin{bmatrix}5&-2\\-1&7\end{bmatrix},\qquad
 b=\begin{bmatrix}4\\-2\end{bmatrix},\qquad c=13.
\]
Indeed,
\begin{align*}
 x^TAx&=\begin{bmatrix}x_1&x_2\end{bmatrix}
 \begin{bmatrix}5x_1-2x_2\\-x_1+7x_2\end{bmatrix}\\
 &=5x_1^2-2x_1x_2-x_2x_1+7x_2^2
 =5x_1^2-3x_1x_2+7x_2^2,
\end{align*}
and $b^Tx=4x_1-2x_2$. Both $x^TAx$ and $b^Tx$ are scalars.

\clearpage
\subsubsection*{Choosing the matrix}
Writing $A=(a_{ij})$ and expanding gives
\begin{equation}\label{eq:lec10-expansion}
 x^TAx=\sum_{i=1}^n\sum_{j=1}^n a_{ij}x_ix_j
 =\sum_{i=1}^n a_{ii}x_i^2+\sum_{i<j}(a_{ij}+a_{ji})x_ix_j.
\end{equation}
Consequently, the construction rule is:
\begin{enumerate}
 \item Put the coefficient of $x_i^2$ in the diagonal entry $a_{ii}$.
 \item Split the coefficient of $x_ix_j$ between $a_{ij}$ and $a_{ji}$, so their sum is the required coefficient.
 \item Put the linear coefficients into the column vector $b$, and keep the constant as $c$.
\end{enumerate}
The off-diagonal split is not unique. In Example~\ref{ex:lec10-two}, all matrices
\[
 A(t)=\begin{bmatrix}5&t\\-3-t&7\end{bmatrix},\qquad t\in\mathbb R,
\]
represent the same quadratic part. Choosing equal halves gives the symmetric choice
\[
 A_s=\begin{bmatrix}5&-3/2\\-3/2&7\end{bmatrix}.
\]
\textbf{A cross coefficient is split in half, not copied into both entries.} Copying $-3$ into both off-diagonal positions would produce $-6x_1x_2$.

\begin{example}[Three-variable construction]\label{ex:lec10-three}
\[
\begin{aligned}
 f(x_1,x_2,x_3)={}&5x_1^2-3x_2^2+2x_3^2
       +16x_1x_2-17x_2x_3-8x_3x_1\\
       &+2x_1-3x_2-4x_3+7.
\end{aligned}
\]
\end{example}
With $x=(x_1,x_2,x_3)^T$, the final symmetric construction on the board is
\[
 A_s=\begin{bmatrix}
 5&8&-4\\8&-3&-17/2\\-4&-17/2&2
 \end{bmatrix},\qquad
 b=\begin{bmatrix}2\\-3\\-4\end{bmatrix},\qquad c=7.
\]
The cross coefficients check as $8+8=16$, $-17/2-17/2=-17$, and $-4-4=-8$. The earlier split $a_{12}=16$, $a_{21}=0$ also works; equal halves simply make the matrix symmetric.

\subsubsection*{The single-variable analogy}
For $f(x)=ax^2+bx+c$ with $a\ne0$,
\[
 f'(x)=2ax+b=0\quad\Longrightarrow\quad x^*=-\frac{b}{2a}.
\]
Since $f''(x)=2a$, this is a minimum if $a>0$ and a maximum if $a<0$. The \emph{optimiser} is the input $x^*$; the \emph{optimal value} is $f(x^*)=c-b^2/(4a)$. Finding a root instead solves $f(x)=0$.

In several variables we cannot divide by a matrix. The analogous product is $A^{-1}b$: an $n\times n$ matrix times an $n\times1$ vector. But we must differentiate to find out which matrix to invert.

\clearpage
\subsection{Gradient of the quadratic form}
\subsubsection*{The linear term}
The coordinates are independent, so
\[
 \frac{\partial x_i}{\partial x_j}=\delta_{ij},\qquad
 \delta_{ij}=\begin{cases}1,&i=j,\\0,&i\ne j.\end{cases}
\]
Here $\delta_{ij}$ is the Kronecker delta. Since $b$ is constant,
\[
 \frac{\partial}{\partial x_j}(b^Tx)
 =\frac{\partial}{\partial x_j}\sum_{i=1}^n b_ix_i
 =\sum_{i=1}^n b_i\delta_{ij}=b_j.
\]
Thus, collecting the partial derivatives into a column vector,
\[
 \boxed{\nabla(b^Tx)=b,\qquad \nabla c=0.}
\]

\subsubsection*{The quadratic term}
Let $q(x)=x^TAx$. In two variables,
\[
 q(x)=a_{11}x_1x_1+a_{12}x_1x_2+a_{21}x_2x_1+a_{22}x_2x_2.
\]
Hence
\[
 \frac{\partial q}{\partial x_2}=a_{12}x_1+a_{21}x_1+2a_{22}x_2.
\]
With more variables, include \emph{every} term containing $x_2$, such as $a_{23}x_2x_3$ and $a_{32}x_3x_2$. Grouping the derivatives gives
\begin{align*}
 \frac{\partial q}{\partial x_2}
 ={}&\underbrace{a_{21}x_1+a_{22}x_2+a_{23}x_3+\cdots}_{(Ax)_2}\\
 &+\underbrace{a_{12}x_1+a_{22}x_2+a_{32}x_3+\cdots}_{(A^Tx)_2}.
\end{align*}
The two copies of $a_{22}x_2$ give the derivative of $a_{22}x_2^2$. Each line is the \emph{second component} of a matrix-vector product, not the sum of all components of that vector.

The same reasoning works for each coordinate $k$:
\[
 \frac{\partial q}{\partial x_k}=(Ax)_k+(A^Tx)_k.
\]
Therefore $\nabla q=Ax+A^Tx=(A+A^T)x$. Adding the linear and constant derivatives gives the result below.

\begin{theorem}[Gradient of a quadratic]\label{thm:lec10-gradient}
For constant $A$, $b$ and $c$, if $f(x)=x^TAx+b^Tx+c$, then
\begin{equation}\label{eq:lec10-gradient}
 \boxed{\nabla f(x)=(A+A^T)x+b.}
\end{equation}
If $A$ is symmetric, this simplifies to $\nabla f(x)=2Ax+b$.
\end{theorem}
\noindent\textit{Study focus:} The lecturer emphasised using the final formulas for the exam. The derivation explains why $A+A^T$ appears, rather than $2A$ for an arbitrary matrix.

\clearpage
\subsection{Stationary point and classification}
\subsubsection*{Why the symmetric part appears}
Recall the symmetric part of a matrix:
\[
 A_s=\frac{A+A^T}{2},\qquad A+A^T=2A_s.
\]
At an unconstrained optimum, the gradient must be zero. Thus
\[
 \nabla f(x^*)=0
 \quad\Longleftrightarrow\quad 2A_sx^*+b=0
 \quad\Longleftrightarrow\quad A_sx^*=-\frac12b.
\]
\begin{theorem}[Stationary point of a quadratic]\label{thm:lec10-stationary}
If $A_s$ is invertible, the unique stationary point is
\begin{equation}\label{eq:lec10-stationary}
 \boxed{x^*=-\frac12A_s^{-1}b.}
\end{equation}
When $A$ is chosen symmetric from the start, $A_s=A$, so $x^*=-\tfrac12A^{-1}b$.
\end{theorem}
Multiply $A_sx^*=-b/2$ on the \emph{left} by $A_s^{-1}$ to obtain this formula. This is why choosing a symmetric matrix during construction is convenient: no later symmetrisation is needed. The constant $c$ affects the function value, but not the location of the stationary point.

\subsubsection*{Classifying the point}
Differentiating the gradient once more gives
\[
 \boxed{\nabla^2f(x)=A+A^T=2A_s.}
\]
The Hessian is constant: there is no $x^*$ to substitute into it. Multiplication by $2$ preserves definiteness, so we classify using $A_s$.
\begin{center}
\renewcommand{\arraystretch}{1.35}
\begin{tabular}{lll}
\toprule
Condition on $A_s$ & Eigenvalues & Conclusion\\
\midrule
Positive definite & All positive & Unique global minimum\\
Negative definite & All negative & Unique global maximum\\
\bottomrule
\end{tabular}
\end{center}
The eigenvalues are real because $A_s$ is symmetric. Positive/negative definiteness concerns the eigenvalue signs, not the signs of individual entries. The two-variable determinant test from earlier lectures does not extend to higher dimensions by checking the determinant alone.

\begin{remark}[Conditions to remember]
The inverse formula requires $A_s$ to be invertible. Setting the gradient to zero finds a stationary point; the definiteness check establishes a minimum or maximum. If $A_s$ is indefinite, a stationary point is a saddle. If $A_s$ is singular, the inverse formula does not apply.
\end{remark}